\renewcommand{\Anno}{2021}

\renewcommand{\NumeroProblema}{2}

\renewcommand{\Testo}{
Per ogni intero positivo $n$ definiamo
\begin{itemize}

\item  $S(n)$ come la somma delle cifre di $n$,

\item  $T(n)$ come la cifra che si ottiene calcolando $S(n)$, $S(S(n))$, $S(S(S(n)))$, e così via fino a quando si ottiene un numero di una cifra sola,

\item  $P(n)$ come il numero di passaggi necessario per calcolare $T(n)$ a partire da $n$.

\end{itemize}

Così, per esempio, $S(76)=7+6=13$ e $S(S(76))=S(13)=1+3=4$, dunque $T(76)=4$ e $P(76)=2$.

\begin{enumerate}
\renewcommand{\labelenumi}{(\alph{enumi})}

\item  Determinare il massimo valore possibile di $P(n)$ al variare di $n$ tra 2.000 e 3.000.

%\item  Determinare il più piccolo intero positivo tale che $P(n)=4$.

\item  Determinare il più grande intero $N$ tale che $P(n)\leq 3$ per ogni intero positivo $n\leq N$.

\item  Scelto a caso (con probabilità uniforme) un intero $n$ tra 1 e 20.000, estremi compresi, determinare i valori più probabili per $T(n)$.

\end{enumerate}

\bigskip
}

\renewcommand{\Risposta}{}

\renewcommand{\Aiutino}{}

\renewcommand{\SoluzioneUfficiale}{}

\renewcommand{\SoluzioniAlternative}{}

\renewcommand{\Osservazioni}{}

\renewcommand{\Commenti}{}
